\documentclass[11pt]{article}
\usepackage[utf8]{inputenc}
\usepackage[T1]{fontenc}
\usepackage{mathpazo}
\usepackage[margin=1in]{geometry}
\usepackage{graphicx}
\usepackage[dvipsnames]{xcolor}
\usepackage{longtable}
\usepackage[htt]{hyphenat}

\newcommand{\code}[1]{\texttt{#1}}

\title{WHAM-XREX-PFMG474 Test Campaign Report\\ \large 10 s ramp / 5 s flat-top}
\author{Generated from the 2026-09-21 controller $\times$ simulated-Transrex campaign}
\date{\today}

\begin{document}
\maketitle

\section*{Summary}

Eight tests were run against the cross-wired controller/simulator rig
(firmware \code{d74b0653-dirty}, 1750-sample PID log). \textbf{6 of 7 discrete
scenarios passed} and the \textbf{45-shot matrix passed 45/45}. The single
non-pass is one criterion-sensitive point in the gain sweep (see \S 1), not a
control defect. Each test below states what was run, why it passed or failed,
and shows its 4-panel plot.

\section{Gain sweep --- closed-loop PID ($P$ + $I$) --- 5/6 PASS}

This test exercised the closed-loop PID controller against the simulated
Transrex plant by commanding a 40\,kHz step on channel~1 under six gain sets:
three P-only ($K_p$ 0.2 / 0.5 / 1.0, $K_i=K_d=0$) and three PI ($K_p=1.0$,
$K_i$ 10 / 50 / 100, $K_d=0$). The P-only runs \textcolor{ForestGreen}{\textbf{passed}}
because the measured frequency settled at the theoretical pure-P steady-state
value $K_p \cdot 40000/(1+K_p)$ --- i.e.\ the expected steady-state error when
there is no integral action. The PI runs \textcolor{ForestGreen}{\textbf{passed}}
because integral action eliminated that error, all three converging to within
$\sim$550\,Hz of 40\,kHz (40\,236 / 40\,255 / 40\,543\,Hz). The one
\textcolor{BrickRed}{\textbf{FAIL}} --- \code{ki100} --- is a verdict-criterion
artifact rather than a control problem: its final sample landed 43\,Hz outside
the $\pm$500\,Hz tolerance while clearly converged, an artifact of judging by a
single ringing-phase sample.

\begin{center}
\includegraphics[width=0.85\linewidth]{../../shots/gain_sweep_4panel.png}
\end{center}

\section{Overcurrent fault injection (\code{ocp}) --- PASS}

Two channels were gated open-loop (ramp 10\,s / flat 5\,s, 3000\,A) and an
overcurrent fault was injected on channel~1 $\sim$0.5\,s into the flat top. The
test \textcolor{ForestGreen}{\textbf{passed}} because \code{STATE?} reported
\code{OVERCURRENT 1}, the faulted channel hard-disabled immediately, and the
surviving channel~2 kept running at a derated 21\,660\,Hz --- exactly the
specified per-channel overcurrent response (immediate disable for the faulted
channel, a $1/N$ step-down then ramp for the survivors).

\begin{center}
\includegraphics[width=0.85\linewidth]{../../shots/20260921_142918_ocp/grid.png}
\end{center}

\section{Enerpro fault injection (\code{enerpro}) --- PASS}

Same 2-channel setup (10/5, 3000\,A), with an Enerpro fault injected on
channel~1 mid-shot. The test \textcolor{ForestGreen}{\textbf{passed}} because
\code{STATE?} reported \code{ENERPRO 1} and the survivor derated and kept
running (21\,638\,Hz), confirming Enerpro faults use the overcurrent-style
\textit{reduce-and-keep-running} response rather than a Water/Temp full stop ---
validating the 2026-09-18 fault reclassification.

\begin{center}
\includegraphics[width=0.85\linewidth]{../../shots/20260921_142935_enerpro/grid.png}
\end{center}

\section{Water/Temp fault injection (\code{watertemp}) --- PASS}

Same 2-channel setup (10/5, 3000\,A), with a combined Water+Temp fault injected
on channel~1 mid-shot. The test \textcolor{ForestGreen}{\textbf{passed}} because
\code{STATE?} reported \code{GENERAL} --- the distinct full-stop fault type ---
confirming that Water/Temp trips the system-wide open-loop ramp-down to 0\,A on
every active channel, rather than the derate-and-survive behavior of OCP/Enerpro.

\begin{center}
\includegraphics[width=0.85\linewidth]{../../shots/20260921_142952_watertemp/grid.png}
\end{center}

\section{ENA\_OUT/CONTACT\_OUT drop (\code{ena\_contact}) --- PASS}

Two channels were gated (10/5, 3000\,A) and the controller's own
\code{CONTACT\_OUT} for channel~1 was dropped mid-shot. The test
\textcolor{ForestGreen}{\textbf{passed}} because \code{STATE?} reported
\code{ENABLE\_OUTPUT 1} and the survivor derated to 21\,660\,Hz --- validating
the per-channel self-commanded-output fault (\code{SM\_FAULT\_ENABLE\_OUTPUT}),
which is distinct from the external-enable \textit{input} interlock.

\begin{center}
\includegraphics[width=0.85\linewidth]{../../shots/20260921_143009_ena_contact/grid.png}
\end{center}

\section{Fault sequencing (\code{multi\_fault}) --- PASS}

Three channels were gated (10/5, 3000\,A) and two different faults were injected
in sequence --- OCP on channel~1, then Enerpro on channel~3 --- with channel~2 as
the survivor. The test \textcolor{ForestGreen}{\textbf{passed}} because both
faulted channels ended disabled, \code{STATE?} preserved the \textit{first}
fault's identity (\code{OVERCURRENT 1}) even after the second fault, and the
survivor visibly derated after each fault (29\,989\,Hz $\to$ 24\,769\,Hz) ---
confirming the ``already-faulted'' short-circuit still hard-disables the
newly-faulted channel without clobbering the reported fault type.

\begin{center}
\includegraphics[width=0.85\linewidth]{../../shots/20260921_143017_multi_fault/grid.png}
\end{center}

\section{Four channels simultaneously (\code{all\_channels}) --- PASS}

All four channels were gated and fired at once (10/5) at per-channel demands of
600 / 2000 / 3600 / 5400\,A. The test \textcolor{ForestGreen}{\textbf{passed}}
because every channel's flat-top median \code{measured~$-$~output} error stayed
well under the 500\,Hz threshold (377 / 387 / 362 / 363\,Hz), confirming four
independent Transrex channels run and measure simultaneously with no
cross-channel interference.

\begin{center}
\includegraphics[width=0.85\linewidth]{../../shots/20260921_143033_all_channels/combined.png}
\end{center}

\section{Channel/demand matrix sweep --- 45/45 PASS}

The stress sweep ran 45 open-loop shots (10/5) --- one per non-empty subset of
$\{1,2,3,4\}$ at each of three demand levels (600 / 3000 / 5400\,A = 10/50/90\%).
Every one of the 45 shots \textcolor{ForestGreen}{\textbf{passed}} with a worst
\code{output~$-$~setpoint} error of \textbf{0\,Hz}, confirming the profile
generator, per-channel HRTIM write, and the fiber feedback round-trip are clean
across the full channel-combination $\times$ operating-point envelope with no
cross-channel coupling.

\begin{center}
\begin{longtable}{l r r l}
\hline
\textbf{Channels} & \textbf{Demand (A)} & \textbf{Worst error (Hz)} & \textbf{Verdict} \\
\hline
\endfirsthead
\hline
\textbf{Channels} & \textbf{Demand (A)} & \textbf{Worst error (Hz)} & \textbf{Verdict} \\
\hline
\endhead
1 & 600 & 0 & \textcolor{ForestGreen}{\textbf{PASS}} \\
1 & 3000 & 0 & \textcolor{ForestGreen}{\textbf{PASS}} \\
1 & 5400 & 0 & \textcolor{ForestGreen}{\textbf{PASS}} \\
2 & 600 & 0 & \textcolor{ForestGreen}{\textbf{PASS}} \\
2 & 3000 & 0 & \textcolor{ForestGreen}{\textbf{PASS}} \\
2 & 5400 & 0 & \textcolor{ForestGreen}{\textbf{PASS}} \\
3 & 600 & 0 & \textcolor{ForestGreen}{\textbf{PASS}} \\
3 & 3000 & 0 & \textcolor{ForestGreen}{\textbf{PASS}} \\
3 & 5400 & 0 & \textcolor{ForestGreen}{\textbf{PASS}} \\
4 & 600 & 0 & \textcolor{ForestGreen}{\textbf{PASS}} \\
4 & 3000 & 0 & \textcolor{ForestGreen}{\textbf{PASS}} \\
4 & 5400 & 0 & \textcolor{ForestGreen}{\textbf{PASS}} \\
1+2 & 600 & 0 & \textcolor{ForestGreen}{\textbf{PASS}} \\
1+2 & 3000 & 0 & \textcolor{ForestGreen}{\textbf{PASS}} \\
1+2 & 5400 & 0 & \textcolor{ForestGreen}{\textbf{PASS}} \\
1+3 & 600 & 0 & \textcolor{ForestGreen}{\textbf{PASS}} \\
1+3 & 3000 & 0 & \textcolor{ForestGreen}{\textbf{PASS}} \\
1+3 & 5400 & 0 & \textcolor{ForestGreen}{\textbf{PASS}} \\
1+4 & 600 & 0 & \textcolor{ForestGreen}{\textbf{PASS}} \\
1+4 & 3000 & 0 & \textcolor{ForestGreen}{\textbf{PASS}} \\
1+4 & 5400 & 0 & \textcolor{ForestGreen}{\textbf{PASS}} \\
2+3 & 600 & 0 & \textcolor{ForestGreen}{\textbf{PASS}} \\
2+3 & 3000 & 0 & \textcolor{ForestGreen}{\textbf{PASS}} \\
2+3 & 5400 & 0 & \textcolor{ForestGreen}{\textbf{PASS}} \\
2+4 & 600 & 0 & \textcolor{ForestGreen}{\textbf{PASS}} \\
2+4 & 3000 & 0 & \textcolor{ForestGreen}{\textbf{PASS}} \\
2+4 & 5400 & 0 & \textcolor{ForestGreen}{\textbf{PASS}} \\
3+4 & 600 & 0 & \textcolor{ForestGreen}{\textbf{PASS}} \\
3+4 & 3000 & 0 & \textcolor{ForestGreen}{\textbf{PASS}} \\
3+4 & 5400 & 0 & \textcolor{ForestGreen}{\textbf{PASS}} \\
1+2+3 & 600 & 0 & \textcolor{ForestGreen}{\textbf{PASS}} \\
1+2+3 & 3000 & 0 & \textcolor{ForestGreen}{\textbf{PASS}} \\
1+2+3 & 5400 & 0 & \textcolor{ForestGreen}{\textbf{PASS}} \\
1+2+4 & 600 & 0 & \textcolor{ForestGreen}{\textbf{PASS}} \\
1+2+4 & 3000 & 0 & \textcolor{ForestGreen}{\textbf{PASS}} \\
1+2+4 & 5400 & 0 & \textcolor{ForestGreen}{\textbf{PASS}} \\
1+3+4 & 600 & 0 & \textcolor{ForestGreen}{\textbf{PASS}} \\
1+3+4 & 3000 & 0 & \textcolor{ForestGreen}{\textbf{PASS}} \\
1+3+4 & 5400 & 0 & \textcolor{ForestGreen}{\textbf{PASS}} \\
2+3+4 & 600 & 0 & \textcolor{ForestGreen}{\textbf{PASS}} \\
2+3+4 & 3000 & 0 & \textcolor{ForestGreen}{\textbf{PASS}} \\
2+3+4 & 5400 & 0 & \textcolor{ForestGreen}{\textbf{PASS}} \\
1+2+3+4 & 600 & 0 & \textcolor{ForestGreen}{\textbf{PASS}} \\
1+2+3+4 & 3000 & 0 & \textcolor{ForestGreen}{\textbf{PASS}} \\
1+2+3+4 & 5400 & 0 & \textcolor{ForestGreen}{\textbf{PASS}} \\
\hline
\end{longtable}
\end{center}

Representative plot (all four channels at the 90\,\% demand level):

\begin{center}
\includegraphics[width=0.85\linewidth]{../../shots/20260921_145413_matrix_ch1+2+3+4_5400A/grid.png}
\end{center}

\section{Closed-loop profiled shot --- PASS (smoke test)}

A one-shot closed-loop counterpart to the profiled shot: channel~1 gated through
the simulator, the PID in \textbf{closed-loop} mode with PI gains ($K_p=1.0$,
$K_i=20$, $K_d=0$), and the same 10\,s/5\,s profile commanded (3000\,A). Unlike
the open-loop runs above, the PID must actually \textit{drive} the output so the
measured feedback tracks the profile. It
\textcolor{ForestGreen}{\textbf{passed}}: the flat-top median
\code{measured~$-$~setpoint} error was \textbf{393\,Hz} (1315 samples @ 52\,Hz,
full 25\,s shot) --- the loop holds the feedback to within the simulator's own
ripple noise of the commanded profile. (The 67\,218\,Hz maximum is the known
single-sample feedback glitch, which the median is robust to.)

\begin{center}
\includegraphics[width=0.85\linewidth]{../../shots/20260921_151030_closed_loop_profiled/ctrl_ch1.png}
\end{center}

\section*{Notes}

\begin{itemize}
\item The \code{gain\_sweep} \code{ki100} FAIL is the final-sample criterion,
  not a control defect (\S 1); a median-of-last-0.5\,s verdict would pass it.
\item Fault tests plot the injection instant as a dashed red line + label on
  every populated panel.
\item Raw per-channel CSVs and all plots are under \code{shots/20260921\_14*\_*/}.
\end{itemize}

\end{document}
